\chapter{Riesgos y Limitaciones del Subsistema SAI}
\label{ch:riesgos}

Este capítulo identifica los principales riesgos técnicos y operacionales
asociados al \textbf{Subsystema de Super--Resolución basada en Inteligencia
Artificial (SAI)}, así como sus limitaciones intrínsecas. El análisis se realiza
siguiendo un enfoque conservador, acorde con el carácter experimental y
complementario del subsistema dentro del proyecto INTISAT.

% -------------------------------------------------------------------
\section{Enfoque de análisis de riesgos}

Los riesgos se clasifican considerando:
\begin{itemize}
    \item impacto potencial sobre la interpretación de resultados,
    \item probabilidad de ocurrencia,
    \item capacidad de mitigación mediante diseño u operación.
\end{itemize}

El subsistema SAI no es crítico para la operación nominal del CMOS Microscopy
Payload, por lo que los riesgos asociados no comprometen la adquisición primaria
de datos.

% -------------------------------------------------------------------
\section{Riesgos técnicos}

\subsection{Generación de artefactos por el modelo de IA}

\textbf{Descripción:}  
Los modelos de super--resolución pueden introducir texturas o patrones
artificiales no presentes en la imagen original.

\textbf{Impacto:}  
Interpretación visual errónea de estructuras microscópicas.

\textbf{Mitigación:}
\begin{itemize}
    \item Uso de parámetros conservadores durante la inferencia.
    \item Comparación sistemática con la imagen original.
    \item Prohibición explícita de uso metrológico de la imagen procesada.
\end{itemize}

% -------------------------------------------------------------------
\subsection{Dependencia del modelo entrenado}

\textbf{Descripción:}  
El desempeño del SAI depende del modelo de IA seleccionado y de los datos usados
en su entrenamiento.

\textbf{Impacto:}  
Resultados variables entre ejecuciones con modelos distintos.

\textbf{Mitigación:}
\begin{itemize}
    \item Registro del modelo utilizado en los metadatos.
    \item Evaluación comparativa entre distintos modelos.
    \item Uso del SAI únicamente como herramienta de apoyo visual.
\end{itemize}

% -------------------------------------------------------------------
\subsection{Limitaciones de generalización}

\textbf{Descripción:}  
El modelo puede no generalizar correctamente a muestras con características muy
distintas a las usadas en entrenamiento.

\textbf{Impacto:}  
Degradación del realce visual o introducción de artefactos.

\textbf{Mitigación:}
\begin{itemize}
    \item Aplicación del SAI solo a imágenes compatibles con el dominio de
    entrenamiento.
    \item Validación visual previa a cualquier análisis.
\end{itemize}

% -------------------------------------------------------------------
\section{Riesgos operacionales}

\subsection{Interpretación indebida de resultados}

\textbf{Descripción:}  
Existe el riesgo de que las imágenes procesadas se interpreten como datos
científicos primarios.

\textbf{Impacto:}  
Conclusiones científicas incorrectas.

\textbf{Mitigación:}
\begin{itemize}
    \item Definición explícita del rol del SAI en la documentación del sistema.
    \item Separación clara entre datos originales y productos derivados.
\end{itemize}

% -------------------------------------------------------------------
\subsection{Dependencia de recursos computacionales}

\textbf{Descripción:}  
El procesamiento de super--resolución requiere recursos computacionales
significativos.

\textbf{Impacto:}  
Limitación en la velocidad de procesamiento o necesidad de ejecución en tierra.

\textbf{Mitigación:}
\begin{itemize}
    \item Ejecución off--line del pipeline.
    \item No dependencia del SAI para operación nominal del payload.
\end{itemize}

% -------------------------------------------------------------------
\section{Limitaciones explícitas del subsistema}

Se establecen las siguientes limitaciones explícitas del SAI:

\begin{itemize}
    \item El subsistema no incrementa la resolución óptica real del sistema.
    \item Los resultados no son válidos para mediciones cuantitativas.
    \item El procesamiento no sustituye técnicas de calibración física.
    \item El SAI no opera de manera autónoma a bordo del satélite en la fase
    actual del proyecto.
\end{itemize}

Estas limitaciones son consistentes con el carácter exploratorio del subsistema.

% -------------------------------------------------------------------
\section{Resumen}

El análisis de riesgos y limitaciones demuestra que el subsistema SAI presenta
riesgos controlables y aceptables dentro del contexto del proyecto INTISAT. La
formulación conservadora del alcance y uso del SAI garantiza que su integración
no compromete la validez científica ni la operación nominal del CMOS Microscopy
Payload.
